\documentclass[../../main-detail.tex]{subfiles}
\begin{document}
\chapter{上位分類}
この章では，コノメノのオントロジー部門で用いる方法論と実際の上位分類の構築について述べる．以下のような構成で進める．
\begin{enumerate}
    \item \emph{方法論}：コノメノではどのような方法論で分類を行うのかを述べる．
    \item \emph{実際の分類}：具体 vs 抽象，持続 vs 生起，独立 vs 従属の既存議論を整理する．さらに，可能世界と虚構物の扱い方を整理し，それを用いて実際の上位分類と，時空間モジュール，量・性質モジュール，表現物モジュールについて述べる．
\end{enumerate}
\section{構築方法論}
\subsection{設計規範}
\paragraph{ハブの使用} ハブ$\langle H, \{\pi_i\}_i\rangle$は，クラス$H$と有限個の二項関係$\pi_i\subseteq V_i\times H$の組である．$\pi_i$は\emph{射影}，$\{\pi_i\}_i$は\emph{射影シグネチャ}と呼ばれる．ハブは開/閉ハブに分けられる．\emph{閉ハブ}は各射影の外延的一致で同一性が与えられる．正確には，関数に変換した射影$\pi_i^\sharp : H \to \wp(V_i)$が jointly injective
（$\forall g,h\in H\ \forall i\;\pi_i(g)=\pi_i(h)\to g=h$）
であるハブだと定義する．閉ハブは，特別な場合として以下を含む：
\begin{enumerate}
    \item $\pi_i$が$H$側について全域かつ一意（つまり$\pi_i : H\to V_i$で関数になる）であるとき，閉ハブに属する要素はタプルと見做せる．
    \item 射影シグネチャに射影を1つしか持たないとき，閉ハブに属する要素は集合と見做せる．
\end{enumerate}
\emph{開ハブ}はそうでないハブである．一般の(開)ハブは$H\twoheadrightarrow R\rightarrowtail \prod_i \wp(V_i)$によって，閉ハブ$R$を介して分解される．Neo-Davidson派（出来事は原始的個物で主題役割はその述語）が開ハブの典型例である．後から外側に射影（時間，様相，記述，評価など）を付け足せることが，開ハブの利点である．

開ハブを標準とし，閉ハブを特殊例とする．閉ハブは形式的構築によって得られる場合にのみ採用し，それ以外の多くの場合は開ハブを採用する．
形式的構築の典型例は\ko{stcilkant}のような純粋な形式的対象であるが，現実世界の対象から形式的構築を通して得られる対象についても閉ハブを採用する．例として\ko{canontoli}は閉ハブであるが，\ko{canon}は開ハブである（\ko{kilem}も同じ型の開ハブ．表現章）．

ハブの導入は，直感的な理解のしやすさから推奨する．典型的に，ハブは
\begin{enumerate*}
    \item コンストラクタ分割，
    \item 引き下げ対応，
\end{enumerate*}
によって導入される．

\paragraph{引き下げによるハブの導入} 3項以上の関係は，意味論で導入した引き下げ対応によって，ハブを導入するのが原則である．具体的には以下の事項を採用する．
\begin{enumerate}
    \item \emph{指標相対化の不採用}：時間変化のための時間指標，モダリティや虚構物のための可能世界指標，記述相対化のための記述指標などを意味値に付け足さない．
    \item \emph{閉ハブの指標付加}：指標が必要なとき，射影として内側に入れず，閉ハブを成分に持つ開ハブを外側に立てる．
    \item \emph{時制は開ハブで表す}：個体$n,p$間に成り立つ関係$nRp$が時間変化するとき，引き下げ対応によって$R$をハブ個体の集合$\tilde{R}$に変換し，「ハブ個体が時刻$t$に存在する」と書くことによって「時刻$t$で$nRp$が成り立つ」と表す．
    \item \emph{スロットを一次概念にしない}：フィールド型・スロット型の記述は採用しない．スロットはハブの射影の呼び名である．
    \item \emph{暫定的な許容}：3項関係やその一部と見做せる2変数関数を使って書くことが自然で簡潔な場合も存在する．この際でも2項関係に直すのが良い（例：\ko{canontoli}）が，3項関係に由来することを注記する．採択が不安定な場合は暫定的に3項関係を使うことも許容する（例：\ko{lop}）．
\end{enumerate}

\paragraph{コンストラクタ分割によるハブの導入} 帰納的定義に似た分類基準がしばしば現れる．これを典型的な分類基準として採用し，コンストラクタ分割と呼ぶ．コンストラクタ分割は，クラス$C$をハブ$H_c$の直和として分類することに当たる．

\begin{leftlinebox}
正規表現($\Sigma$を文字集合とする)
$M ::= \Sigma \mid M M \mid M + M \mid M^*$
をオントロジーで書くと，
\begin{figure}[H]
    \centering
    \begin{forest}
        dir tree
        [(root),draw=none,
            [\tn{$\Sigma$}{文字}, name=letter]
            [\tn{$M$}{正規表現}
                [単純表現, name=atomic
                ]
                [複合表現
                    [連接, name=concat]
                    [選択, name=choice]
                    [繰返し, name=repeat]
                ]
            ]
        ]
        \treerel[via=rightup]{concat}{atomic}{\scriptsize $\pi_1$, $\pi_2$}
        \treerel[via=rightup]{choice}{atomic}{\scriptsize $\pi_1$, $\pi_2$}
        \treerel[via=rightup]{repeat}{atomic}{\scriptsize $\pi$}
        \treerel[via=up]{letter}{atomic}{\scriptsize 見做し}
    \end{forest}
    \caption{正規表現のコンストラクタ分割}
\end{figure}
\end{leftlinebox}
具体的には以下の事項を採用する．
\begin{enumerate}
    \item \emph{列挙}：列挙は0項のコンストラクタによる分類としてコンストラクタ分割に含める．
    \item \emph{再帰の許容}：コンストラクタ分割により，高階の対象が無限に生成されることは仕様であり抑制しない．例：\ko{canon}が担い手になれることから高階\ko{canon}が作れ，\sentence{とても長い}は別単語を立てずに書ける．表現物にも同種の再帰がある．
    \item \emph{複数の表示を持つ場合}：複数の帰納的な表示を持つもの（整数の0／正／負と差の同値類など）はコンストラクタ分割だと見做さない．
\end{enumerate}

\paragraph{同一性の様式を宣言する} 語（特にハブと種）を立てるときは，その同一性の様式を宣言する：
\begin{enumerate*}
    \item 原始的（開ハブ．Leibniz則の否定方向，すなわち両立しない述語を持てば別個体という区別の十分条件と，種が供給する計数基準で運用する），
    \item 射影外延的（閉ハブ），
    \item メレオロジー的（状況．部分が同じなら同じ．fusionを持つ），
    \item 構成員外延的（\ko{klant}・集成生起物．成員が同じなら同じ）．
\end{enumerate*}
様式が宣言できない候補はゴミ箱分類群の疑いがあるので攻撃工程に回す（判例：\ko{nosnant}は同一性の様式を先に決めなかったためゴミ箱化した．未決の「日 vs 日付」「色 vs 色の種類」も正体は同一性の問いである）．
\todo{本当にこれに沿ってやっているのか，必要なのか}

\paragraph{要素分解} 役割や語彙概念は原始的に立てず，少数の原始関係——担い手射影，\ko{nila}，\ko{skon}，\ko{moila}／\ko{fiila}，\ko{stae}，因果先行$\mathrm{luis}^{+}$（直接因果\ko{luis}はその被覆関係として定義．2026-09-03），\ko{na}——の組合せとして定義することを目標とする（2026-09-02採択）．定義できないものだけを原始に残し，原始関係の追加には新装置と同じ立証責任を課す．採択判例：\ko{kalastav}廃止（完遂＝\ko{fiila}の存在言明），\ko{tact}廃止（使役＝\ko{luis}接続），\ko{towa}＝中間に独立の内部駆動源を経ない\ko{luis}連鎖，\ko{e}＝正の段階の開始または終了境界が非状態トークンの終端と共有・因果されること，\ko{ja}＝$\mathrm{ja}_{\mathrm{src}}$（\ko{tae}＋意図）と$\mathrm{ja}_{\mathrm{res}}$（\ko{towa}＋意図起動）の型付き和（decisions/0003，2026-09-02）．検証は標準形／分解形の同値ペアをcorpusに置き，双方向の含意と反例で行う．（\ko{astav}の部分関係$\mathrm{incl}$は2026-09-05に原始に数えた．生起物の章「個体在庫と原始関係」）
\todo{短くする}

\paragraph{見做し} 存在論的には別の節点だが，語彙では呼び分けない対がある：\sentence{学校}と\sentence{学校のあるところ}，\sentence{本の内容}と\sentence{物理的な本}，\sentence{物質としての花瓶}と\sentence{オブジェクトとしての花瓶}．これらを結ぶメトニミー的な二項関係の慣用的な呼び方を\emph{見做し}と呼ぶ．見做しは木の辺ではなく，
\begin{enumerate*}
    \item 存在論上の二つの節点，
    \item それを結ぶ指定された二項関係，
    \item 両端を呼び分けないという語彙化の決定，
\end{enumerate*}
の三点セットである．述語の選択制限がどちらの側面を取るかを決める（言語学の規則的多義性・ドット型に当たる）．

見做し関係は有限の目録として与える．抽象/具体の節に挙げた依存関係のうち，存在論的構成（像と粘土），媒体への一般依存（内容と物理的な本），組織と建物と場所の宿り（学校）が該当し，これに帰納的定義の原子の注入（文字を単純表現と見做す）と一部の1項コンストラクタを加える（目録の範囲は，持続物の章「見做しグラフと語義分割」で有限個の型付き二項関係からなるグラフへ一般化した．2026-09-03）．

見做しにするか呼び分けるかの基準は，意味移行の規則性である：その移行がクラス全体に一様に適用できる（すべての学校に所在地があり，すべての本に内容がある）なら見做し，個別的なら呼び分ける．剛性の違いは見做しを疑う徴候にはなるが基準ではない．「関係のドメインには節点を立てる」（分類作法）は節点についての規則であり，見做しはその節点に固有の語を与えない意図的な例外である．
\todo{短くする}
\todo{schedule2から取ってくる}

\paragraph{配置の状態} 各辺は配置の状態として，\emph{配置}ないし\emph{上位未決}のいずれかを宣言する．前者は包摂を主張するが，後者は分かりやすい位置に置いているだけで包摂を主張しない．これは現状のオントロジーで扱いきれない語（対応予定があるものと，他言語にあるがオントロジーでは扱う予定がない語）を扱うための措置である．

\paragraph{分割} ある節点直下の典型的な分類基準（\emph{分割}と呼ぶ）に任意で注釈をつけることができる．分割には以下の2種類がある（全てがこれらに当てはまることは主張しない）．
\begin{enumerate}
    \item \emph{コンストラクタ分割}：枝ごとに射影シグネチャが違う．外延的な列挙（0項のコンストラクタ）とソートの分岐（下の図の文字 / 正規表現）もこれに含む．自動的に排他性が満たされる．例：\ko{nato}の下の\ko{matonakt} / \ko{noflem}，正規表現の単純表現 / 複合表現，\ko{stcilkant}の直下．
    \item \emph{素性分割}：全成員が同じ射影シグネチャを共有し，ある属性の値で分ける．交差してよく，排他かどうかは軸ごとに宣言する．軸と値のクラスにも節点を与える．例：\ko{astav}の変動パタン，\ko{takto}の統一源泉．
\end{enumerate}
次の場合は分割の注釈をつけることを推奨する：
\begin{enumerate}
    \item 複数軸による素性分割があるとき，軸ごとに宣言する．
    \item 分類基準に沿わない例外的な語を配置する．（このような語は旧来は「未配置」と呼ばれ，配置の状態に含まれていたが無用になった．上位未定とは違って包摂を主張するが，分類基準に沿わない）
\end{enumerate}
また，分割には公理が追加することができる：
\begin{enumerate}
    \item \emph{排他}：枝の外延が交わらないことを主張する．
    \item \emph{網羅}：枝の外延の和が節点の外延と一致することを主張する．
\end{enumerate}

\memo{時間とともに値が変わる素性（幼虫／成虫，相）で持続物を分割してよいか．指標を入れない以上，可変な素性は\ko{canon}側に置くことになるはず．ロール（\ko{liflom}）も同じ扱い．}
\memo{確定可能者／確定者（色$\to$赤）は属＋種差で書けない．値空間の領域による分割として扱うか．}

\paragraph{粒度中立性} 分類は特定の粒度にコミットしない．要求するのは，粗い粒度の分類が細かい粒度の分類の商になっていることである（時空間モジュールのパタン型と同じ作法）．粒度には意味的なもの（人工物のロール性）と時空間的なもの（イベントの境界，\citet{tawatsuji2021}）がある．
\memo{実用上どうなのかわからない．例を集めることが必要である}

\paragraph{種を個体として立てる基準}
\begin{enumerate}
    \item 種をまとめる語がある．（例：\sentence{職業}$\sqsupseteq$\sentence{教師}）
    \item 成員かどうかの判断が相対的．（例：\sentence{モンスター}．反剛性）
    \item 種水準の述語が当たる．外延の各成員には当てはまらない述定．（例：\sentence{絶滅した}，\sentence{発明された}，\sentence{広く分布する}）
    \item 空外延・共外延で区別が要る．（例：de dicto目的語のユニコーン$\neq$ペガサス．クラス唯名論への標準批判が当たる箇所）
    \item 「〜という種」への照応・計数がある．（例：\sentence{この属には3種いる}）
    \item 同一性の様式を供給する．（成員の計数基準を与えるか）
\end{enumerate}
\todo{種を個体として立てるのはなんでもありになるのでもっと慎重に考えた方がいい}

\subsection{検証方法}

\paragraph{OntoCleanによる分類の検査}
OntoCleanは，クラスの剛性・同一性・統一性を明示してis-aを検査する方法である（\citet{guarino2009ontoclean}）．以下は検査用のメタ言語であり，$P,Q$はクラス，$E(x)$は存在，$\Box,\Diamond$は必然・可能を表す．$Q\sqsubseteq P$は$\Box\forall x(Qx\to Px)$とする．
\begin{enumerate}
    \item \textbf{剛性（rigidity）}：
    \[
    \begin{aligned}
        +R(P)&\iff\Box\forall x\bigl(Px\to\Box(E(x)\to Px)\bigr),\\
        {\sim}R(P)&\iff\Box\forall x\bigl(Px\to\Diamond(E(x)\land\neg Px)\bigr).
    \end{aligned}
    \]
    $-R$は$+R$の否定であり，反剛性${\sim}R$より弱い．\textbf{テスト}：同じ個体が存在し続けながら$P$でなくなれるか．全成員で可能なら${\sim}R$，一例でも可能なら$+R$を却下する．退職しても同じ人なので教師は${\sim}R$．単なる消滅や呼称の変更は反例にならない．

    \item \textbf{同一性（identity）}：$+I$は同一性基準を担うこと，$+O$は上位からの継承でなく自ら供給すること．$P$の成員$x,y$について，候補条件$C$が
    \[
        C(x,y)\to x=y\quad\text{（十分条件）},\qquad
        x=y\to C(x,y)\quad\text{（必要条件）}
    \]
    のどちらを満たすかを分ける．通時的な基準では$C(x,y,t,t')$として比較時点も指定する．$C$を同一性そのものの言い換えにせず，十分条件は充足可能，必要条件は成員の任意の対に自動成立しないものとする．\textbf{テスト}：$C$を共有する別個体があれば十分性を，同じ個体が$C$を失えるなら必要性を却下する．同じ長さの別の時間区間があるので，長さの一致は時間区間の同一性の十分条件でない．成員資格の条件と同一性基準も別である．

    \item \textbf{統一性（unity）}：時点を固定し，対象を尽くす部分の非空集合$B$と，部分間の統合関係$R$を与える．次の二条件を分けて検査する（\citet[定義3--4]{guarino2000identityunity}の条件を分解したもの）．
    \[
    \begin{aligned}
        \operatorname{Con}_R(B)&\iff\forall x,y\in B\;\bigl(Rxy\lor Ryx\bigr),\\
        \operatorname{Cl}_R(B)&\iff\forall x\in B\;\forall y\;\bigl((Rxy\lor Ryx)\to y\in B\bigr).
    \end{aligned}
    \]
    前者は内部の結合，後者は外部に漏れない$R$-閉包性である．基礎関係$r$が直接の接触などなら，$R=(r\cup r^{-1})^*$（反射推移閉包）を用いる．\textbf{テスト}：内部の未結合の対，外部へ結ばれる相手を探す．$R$を「同じ$B$に属する」と逆算してはならない．
    $+U$は全成員が共通の基準で本質的にwholeとなること，$-U$は共通基準を持たないこと，${\sim}U$はどの成員にもwholeであることが本質的でないこと．実際に一塊であるだけでは$+U$を示さない．\suppl{非空集合$B$ごとに$R_B(x,y)\iff(x\in B\land y\in B)\lor x=y$と置けば，同値関係のもとで結合・閉包が自動成立する．したがって関係の形式的性質だけでなく，対象の成員指定から独立した$R$の選択根拠が必要である（\citet[第4.1節]{guarino2000identityunity}の非自明性条件）．}この時点での部分のまとまりは，部品交換後も同じ個体かを決める同一性基準とは別である．
\end{enumerate}

\noindent\textbf{is-aの検査}：$Q\sqsubseteq P$なら，$P$の${\sim}R$，同一性基準，統一性基準，${\sim}U$は$Q$へ継承される．したがって${\sim}R$の下に$+R$，${\sim}U$の下に$+U$を置けず，多重継承では両親の基準を同時に満たす必要がある．たとえば海洋（$+U$）は水の量（${\sim}U$）の下位種でなく，水から構成される．タグの一致だけではis-aを証明できない．ロールは\ko{liflom}としてプレイヤの木の外に置く．

\paragraph{言語使用によるテスト} 曖昧な判断基準（例えば「本質」「依存」など）は容認性テストを用いて明確化する．言語上での容認性と存在論的な正確性は必ずしも一致しないという指摘があるが，他の基準と組み合わせて使用する前提のもとで，
\begin{enumerate*}
    \item 言語制作という目的によく合うため，
    \item 客観的基準として強力であるため，
\end{enumerate*}
採用する．
\begin{enumerate}
    \item \emph{表現可能性}：ある表現が可能かどうか．3つの分類が関係づけられる場合によく見かける．「対象$N$が量$t$を持つことは形容$a$である」「時空領域$A$は系$R$でラベル$v$を持つ」など．
    \item \emph{共述定}：見做しの判定．両側面を一文で述定できれば見做し（\sentence{学校は築50年で厳しい}，\sentence{本は重くて退屈だ}），座りが悪ければ呼び分け候補．
    \item \emph{計数}：同一性の様式の判定．「二つの職を持つ一人は教師として二人か」のように計数が変われば別個体への写像であり，is-aではない．
\end{enumerate}

\paragraph{判定テストと境界例} 分類基準ごとに判定テストを設計する．適当な自然言語の表現で定義した気にならない．言語表現による定義でなくてもよい．うまく定義できないなら外延的に与える（例：\ko{tastoi}/\ko{liskono}，具体物の派生的位置）．感覚が先行するので，境界例を積極的に集めてテストにかける．現在の境界例：小さい象と大きな蟻（閾値），日と日付，色と色の種類（類とインスタンスの二重性），像と粘土（構成），魚（側系統群），チェス（準抽象物），18cmに戻った棒（開\ko{canon}の再発），走り終える／やめる（完了と停止），見られただけの対象（随伴と帰結），描きかけの絵（\ko{e}の事実性と名詞化子の部分性）．

\section{抽象/具体}
参考：\citet{sep-abstract}．抽象/具体の一般的に知られている定義は「時空間上に位置を持つかどうか」である．
\citet{mizoguchi2012}に倣えば，「時間上に位置をもつ」($\Leftrightarrow T$)，「空間上に位置をもつ」($\Leftrightarrow S$)の2つの基準を組み合わせて次の3分類を得る（$S \land \lnot T$は存在しないとされる）：
\begin{enumerate}
    \item $S \land T$：具体物，
    \item $\lnot S \land T$：準抽象物，
    \item $\lnot S \land \lnot T$：抽象物．
\end{enumerate}
DOLCE的に整理するなら以下のようになる：
\begin{enumerate}
    \item $T$：concrete
    \item $S \land T$：spatial concrete
    \item $\lnot S \land T$：non-spacial concrete / quasi-abstract
    \item $\lnot T$：abstract
\end{enumerate}
\todo{たぶん今後よく参照するのでDOLCEの原典は読んだ方がいいかもしれない．}
準抽象物はチェスのような発明された``抽象的な''人工物を想定しているが，抽象物との区別が必ずしも明確ではない．これについてはあとで議論する．次によく用いられる定義が，「因果的効力を持つ($\Leftrightarrow$具体物)かどうか」であるが，（因果についての直感がないため）この定義は避けることにする（そんなことしていいの？）．

時空間への位置付けによる定義をもう少し深掘りする．素朴には「個体$x$が場所$S$にある」と言えるのが具体物である．典型的な持続物（「机」「人間」など）が具体物に分類されることは認めるとし，その上で (BFO・DOLCEなどの一般的なオントロジーでは) 生起物はその参与者，集団はその構成員を通じて空間上に\emph{メレオロジー的派生位置}を持つため具体物に分類されると考える．このように，何かしらの構成要素と判断したい対象を繋ぐ\emph{構成的}関係を通して位置を誘導することが認められる．
参与を「構成的」関係と呼ぶのは用語として不正確なので，ここでは位置を誘導する関係の総称として読む．
\begin{remark}
集団\ko{klant}は具体物，不純集合は\ko{loka}の形式構築物（抽象物）として語彙を分ける．構成員関係\ko{hanem}は派生的位置を与えるが，所属関係\ko{litoi}は与えない．この関係の選択は規約であり，両者を区別する非恣意的な根拠は未決である．
\end{remark}

これに対して，ロゴなどの反復可能なパタンや，「18cm」「130g」などの量，「教師」「朝食」などのロールは抽象物に分類される\suppl{ロールのプレイヤが抽象物だというわけではない}．それは，実行や実例化（さらには所有や因果等も同様）は構成的関係とは見なされないためである．

似ているが異なる機構として，canon や haikantのような従属物（自身の部分だけでなく外部の個物に依存する）は，担い手（bearer）を通じて\emph{依存的派生位置}を持つと考えられ，これも具体物に分類されると考えることが多い．ただし，``依存''と呼べる関係は多岐にわたる（あとで依存性の考察をする際にそちらにまとめたい）：
\begin{enumerate}
    \item bearerへの特定依存：このリンゴの色，この人の痛み
    \item 媒体への一般依存：小説，ロゴ，プログラム
    \item コンテキスト依存：教師ロール，貨幣，所有権
    \item 歴史的依存：作品の作者への依存
    \item 因果的依存：製品の製作者への依存
    \item 存在論的構成依存：像と粘土
\end{enumerate}
したがって，どの種類の依存関係を許容するかは明示的に列挙して与える必要があるだろう．ここまで整理すると，具体物の位置の持ち方は少なくとも，固有位置，メレオロジー的派生位置，依存的派生位置の3つがあることになり，位置を継承・集約する公理を個別に与えられた構成・参与・依存関係等に基づいて「空間的位置を持つ」と定義することができる．
\memo{外延的に基準を与えるのは，明確になる分，良いが恣意的ではないか？}

\section{世界や形式対象} コノメノの従来の具体 / 抽象の分類は「時空間上に位置を持つか」というより「時空間なしに存在できるかどうか」に近い．もっと言えば「この世界が無くなったとしてもAは存在するのか？」(A)という問いに基づくのだが，実際に世界が無くなったことがあるわけでもないし，この問いはそれぞれの人の存在論的なコミットメントに依存してしまう．特に，この基準に基づいて形式的対象を``抽象物''に分類すると，暗黙に（形式的対象の実在を認めるという）プラトニズムを前提としてしまう．これを回避するためには，形式的自律性「概念$P$の同一性基準が形式的体系の内部で完結しており，時空間・因果・感覚経験への参照なしに$P$が何であるかが完全に決定される」(B)という基準で抽象物の基準とした方が良い．

(A)または(B)の基準に基づくと，``具体物''がかなり多くなる（時空間や現実世界の何かしらを参照するものはすべて具体物に分類される）．``抽象物''に属すると想定される典型例は形式的対象である：形式的対象は，「四角形」のように現実に近似的に実現できるものもあるが，それはほんの一部であり，同一性基準が形式体系の内部で完結している．

(B)は基本的に形式的対象の特徴づけであるが，(A)の場合はプラトニズムを認めてしまえば，他の世界（小説内の世界）も同じ基準で扱える．本版は(B)を採る（2026-09-03裁定）：(A)は虚構物を現実の対象と同列に扱えるという利点があるが，現行の他の機構（表現章の虚構物は充足世界クラスと対応者関係で扱い，現実の対象と同列にはならない）を見るとその利点は活かせないので，このバージョンでは(B)が自然である．(A)は将来の版のための代替案として残す．要するに，最上位分類は次のようになる：

\begin{figure}[H]
  \centering
  \begin{forest}
    dir tree
    [概念
      [世界
        [現実世界]
        [$w_1$, dashed]
        [$w_2$, dashed]
        [$\cdots$, dashed]
      ]
      [形式的対象]        
      [$w_1$の個物, dashed]
      [$w_2$の個物, dashed]
      [$\cdots$, dashed]
      [現実の対象]
    ]
  \end{forest}
\end{figure}
世界を異にする対象も，共通の構成骨格と語義の定義条件に従って分類する（2026-09-06採択，暫定）．共通なのは分類上のすべての包摂事実ではない：構成公理と，同じ語義を保つために必要な定義公理を共通公理$\Gamma_{\mathrm{c}}$とし，各世界の語彙解釈$I_w:=\rho\vll\rho_w$はこれを満たすものに限る．経験的な種分類，制度上の所属，個体の分類などは，各世界の外延に対する制約$\Gamma_w$として与える．
\[
\forall w\,[\,I_w\models\Gamma_{\mathrm{c}}\cup\Gamma_w\,],\qquad
\Gamma_{\mathrm{c}}\models A\ \ko{skon}\ B\ \Rightarrow\ \forall w\,[\mathrm{Ext}_w(A)\subseteq\mathrm{Ext}_w(B)],\qquad
I_{w_0}\models A\ \ko{skon}\ B\ \not\Rightarrow\ \forall w\,[\mathrm{Ext}_w(A)\subseteq\mathrm{Ext}_w(B)].
\]
\ko{skon}は集合包含を表し，その意味と推移律は世界によって変わらない．同じ概念について包摂文の真偽が世界によって変わるのは，その概念の外延が変わるためである．異なる世界で成立する包摂文を，世界の指定を落として一つの推論へ混ぜてはならない（中間項の集合が同じだと保証されないからであり，推移律の例外ではない）．「コンストラクタ分割だけが必然で他はすべて偶然」ではなく，定義からの包含も全世界で維持する．逆に，剛性から上位分類の全世界不変性は導けない（剛性は同じ個体の所属に関する制約であり，現実に存在しない可能的実例まで拘束しない）．
世界別の分類と個体の記録は，既存フレームの語彙割り当てによって解釈する．記録を世界ごとの台帳として分けるかは保存上の選択であり，台帳を新しい存在者の種や個体生成の機構とはしない．辞書への未登録は非存在を意味しない．作中の語と現実の語が同じ語形を持つことは同じ概念を表す十分条件ではない：同一性条件と定義条件を保つ場合には同じ概念の世界ごとの外延として扱い，条件を変更する場合には別概念を立てる．対応・類似・翻訳上の呼び換えを\ko{skon}と同一視しない（表現章「虚構物」）．
\begin{remark}[世界際視点]
    現実世界の視点だけでは足りない理由は，主に嘉音もコノメノを話すためである．嘉音は世界際的な存在であり，視点を現実世界に固定することは自明に正当化できない．作中人物の発話を基準とする場合は，フレームの基底解釈をその世界へ再基底化し（$\rho' := I_v$，各$w$について$\rho'\vll\rho'_w = I_w$となる差分$\rho'_w$を取る．したがって$\rho'_v=\emptyset$で，有限差分どうしの再基底化なら差分の有限性も保たれる），発話状況と談話状態（談話モジュール）を当該発話に合わせる．この操作は世界際的な個体同一性や移動（メタレプシス）を決定しない．それらは別途，対応ハブと作品内容によって記述する（2026-09-06，暫定．旧稿は「工数の多さからしばらく無視する」としていた）．
\end{remark}
虚構物は表現章「虚構物」と前段の世界依存の分類で扱った（2026-09-06）．\todo{不純集合（集合化で日常的に出てくる）は(A)でも(B)でも現実の対象に属する．「形式構築物」というカテゴリ（普遍者や方向のような同一視の結果を含む）を立てるかは未決．}

虚構物，ゲーム内オブジェクト，理想化された系，数理モデルの扱いは表現章「虚構物」節に集約した（2026-09-06）：層W／K／Lの図式，役と世界内個体の対応ハブ，解釈済み文への冠詞$K$と背景の持ち込みによる作中帰結，内容の形式化（形式化$\subseteq$\ko{klaleme}$\times$\ko{stcilkant}，一意でない）．本章では前節の世界依存の分類（共通公理と世界ごとの制約）だけを扱う．

\section{実際の分類}

\begin{figure}[H]
  \centering
  \begin{forest}
    dir tree
    [\tn{alkono}{概念}
      [\tn{suis}{世界}]
      [\tn{stcilkant}{形式的対象}]
      [\tn{suiskant}{現実の対象}
        [\tn{toika}{具体物}
          [\tn{nato}{持続物}]
          [\tn{astav}{生起物}]
          [\tn{aicis}{個別従属物}
            [\tn{canon}{トロープ}]
            [\tn{haikant}{特徴}]
            [(個別のロール)]
            [(個別の傾向)]
          ]
        ]
        [\tn{kika}{準抽象物}]
        [\tn{loka}{抽象物}
          [\tn{latent}{量}
            [\tn{lekostent}{定値量}]
            [\tn{nosnant}{形容量}]
          ]
          [\tn{liflom}{ロール}]
          [(評価記述)]
          [(傾向タイプ)]
          [(形式構築物)]
          [(概念個体)]
        ]
        [\tn{metlosk}{基盤}
        ]
      ]
    ]
  \end{forest}
\end{figure}

\section{独立/従属}

\section{持続/生起}

\section{時空間モジュール}
\todo{省略したが kakt vs hoft 数値の議論の置き場所が必要．}
時空間と，対象を時空間へ位置付けるために必要な構造をまとめて
metloskと呼ぶ．この呼び方はYAMATOの「基盤」に倣ったものである．
基盤は，具体物と抽象物の中間物として導入するのではなく，両者の区別を成立させる
上位分類として導入する．

日常的な範囲で，時間と空間の大きな違いは「基準系について相対的($\Leftrightarrow$空間)か絶対的($\Leftrightarrow$時間)か」である．これは言語の上にも現れると考えている：例えば「家の前」「箱の中」のような空間表現は，ふつうは物体を基準にして記述される．それに対して，「2024年1月5日」「朝9時」のような時間表現は，時間の位置をそのまま指定する．日常的な範囲では，ローレンツ変換を実装する必要はないが，どちらにしても基準系について相対的な扱いが必要であるから，両方とも扱えるような仕組みを用意した方が良い．

\subsection{時空間と参照付き位置付与}

個々の時空間を，必要な構造を備えた数学的構造
\begin{align*}
    \mathfrak{S}
    =
    (M,\tau,\mathcal{A},g,\nabla,\mathfrak{o},\ldots)
\end{align*}
と同一視して説明する．ここで$M$は台多様体，$\tau$は位相，$\mathcal{A}$は
微分構造，$g$は必要に応じて入る計量，$\nabla$は接続，$\mathfrak{o}$は向きなどを表す．
どの成分を要求するかは対象領域によって異なるため，以下では構造全体を単に$M$と略記する．

この同一視では，次の三つを区別する．
\begin{enumerate}
    \item 構造付き集合として文字どおり等しいこと，
    \item 微分同相・等長写像など，指定された種類の同型で移り合うこと，
    \item 同じ物理的時空間を表すとみなすこと．
\end{enumerate}
時空間の型の同一性は(2)によって与える．一方，個々の点や領域を指示する場面では，
必要に応じて$(\mathfrak{S},p)$のような点付き構造を用いる．

一般の時空間$M$には，参照構造から独立した分解
\begin{align*}
    M\simeq T\times S
\end{align*}
が備わっているとは限らない．したがって，基盤の最初の層では時間的対象と空間的対象を
分けず，時空点または時空領域$A$を扱う．時間位置と空間位置への分解は，参照構造$R$を
与えた後に初めて得られるものとする．

時空間の位置記述には，次の三つを区別する．
\begin{itemize}
    \item $A$：参照構造から独立に取られた時空点・時空領域・そのパタン，
    \item $R$：$A$を観測・分解・表示するための参照構造，
    \item $v$：表現体系によって型付けられているが，特定の$R$にはまだ束縛されていない値．
\end{itemize}
例えば，「午前2時」という値は12時間制の時計表現には依存するが，それだけではJSTや
特定の日付を指さない．「北緯35度」も緯度という値形式には属するが，測地系が与えられる
までは一つの時空領域を指示しない．この意味で，時空間値は「参照系から独立」なのではなく，
より正確には「特定の参照構造に未束縛」である．

三項関係「$A$は$R$のもとで$v$を持つ」を二項関係だけで扱うため，これを仕様化したハブ
canontoliを置く．その個体$h$は
\begin{align*}
    h=\langle A,R,v\rangle,
    \qquad
    v\in V_{\operatorname{scheme}(R)}
\end{align*}
と表され，三つの射影
\begin{align*}
    \pi_A(h)=A,\qquad
    \pi_R(h)=R,\qquad
    \pi_v(h)=v
\end{align*}
を持つ．単純な位置付与としての同一性基準は
\begin{align*}
    h=h'
    \quad\Longleftrightarrow\quad
    \pi_A(h)=\pi_A(h')\land
    \pi_R(h)=\pi_R(h')\land
    \pi_v(h)=\pi_v(h')
\end{align*}
とする．同じ$A,R,v$についても，情報源・作成者・精度・記録時刻の異なる記述を区別したい
場合，それはcanontoliそのものではなく，canontoliを内容として持つ情報対象として
kika以下に置く．また，紙面や音声に実現した記述トークンは具体物である．

\begin{figure}[H]
  \centering
  \begin{forest}
    dir tree
    [\tn{metlosk}{基盤}
      [\tn{losktoli}{時空間関連}]
      [\tn{losktseltoli}{時空間値関連}
        [\tn{tseltoli}{時空間値}]
        [\tn{tselselov}{時間値}]
        [\tn{tselfenov}{空間値}]
      ]
      [\tn{notlileka}{参照構造}]
      [\tn{canontoli}{参照付き時空位置}]
    ]
  \end{forest}
  \caption{基盤の主要分類}
\end{figure}

losktoliは，参照構造から独立した$A$の側を分類するため，時間と空間を分離しない．
これに対して，losktseltoliは$v$の側を分類する．時間値と空間値の区別は，$A$そのものの
分割ではなく，$R$によって定まる射影のどちらの余域に値が属するかを表す．

\subsection{点・領域・複合パタン}

以下では，時空間$M$を少なくとも位相多様体とし，必要に応じて微分構造・計量・群作用を
追加する．多様体であることだけから距離・長さ・合同変換が得られるわけではないため，
それらを用いる箇所では追加構造を明示する．

$M$上の点・領域・それらの族のうち，以下の定義によって位置付けの対象として認めるものを
まとめた領域を$\mathfrak{L}(M)$と書く．これは$M$の冪集合そのものではなく，点，正則閉領域，
局所有限族などを型付きで合わせた位置対象の領域である．

離散的な点の集まりには，一様離散部分集合ではなく局所有限部分集合を用いる．
\begin{align*}
    \mathcal{D}(M)
    =
    \left\{
      P\subseteq M
      :
      \forall p\in M\;\exists U\ni p,\; |U\cap P|<\infty
    \right\}.
\end{align*}
これにより有限集合や無限格子を含めつつ，有限位置に集積する点列を除外できる．
$M$が局所コンパクトな場合には，各コンパクト集合$K\subseteq M$に対して
$|K\cap P|<\infty$であることと同値になる．

図形的な領域には，連結部分多様体よりも正則閉集合を用いる．
\begin{align*}
    \operatorname{RC}(M)
    &=
    \{H\subseteq M\mid H=\overline{\operatorname{int}H}\},\\
    \mathcal{H}(M)
    &=
    \{H\in\operatorname{RC}(M)\mid H\text{は連結}\}.
\end{align*}
これにより，境界や穴を持つ領域も扱える．非連結領域は，連結領域の族として表す．
領域族$\mathscr{H}\subseteq\mathcal{H}(M)$の局所有限性は
\begin{align*}
    \forall K\subseteq M\text{ compact},\qquad
    \#\{H\in\mathscr{H}\mid H\cap K\neq\emptyset\}<\infty
\end{align*}
によって定め，そのような領域族のクラスを$\mathcal{D}_M(\mathcal{H}(M))$と書く．

単一・複合・繰り返し・恣意的という分類は，点・領域・長さを横断する生成様式である．
基本形は次のように表す．
\begin{figure}[H]
  \centering
  \begin{forest}
    dir tree
    [$x$-パタン
      [単一$x$]
      [複合$x$
        [繰り返し$x$]
        [恣意的$x$]
      ]
    ]
  \end{forest}
  \caption{パタンの生成様式}
\end{figure}
ここで「複合」は複数の実現を持つことを表し，「繰り返し」はその複合が一定の生成規則を
持つ場合，「恣意的」は共通の生成規則を指定しない場合を表す．したがって，繰り返しを
時点・区間・長さと並ぶ第四の対象種とはしない．

群$G$が$M$または領域族に作用するとき，$X$の繰り返し実現は軌道
\begin{align*}
    G\cdot X=\{gX\mid g\in G\}
\end{align*}
によって与えられる．これに対して，$[X]_G$は$G$で移り合う対象を同一視したパタン型であり，
商空間$\mathcal{X}/G$はそのような型全体の空間である．したがって，次の三つを区別する．
\begin{align*}
    &\text{繰り返し実現} &&G\cdot X,\\
    &\text{一つのパタン型} &&[X]_G,\\
    &\text{パタン型の空間} &&\mathcal{X}/G.
\end{align*}
特に，繰り返し実現そのものを$M/G$と書かない．また，群軌道が局所有限になるためには，
$G$が離散的で作用が適切に不連続であることなど，追加の条件が必要である．連続回転群の
点軌道は円をなし，離散的な点パタンではなく領域的な対象となる．

長さ・広がり・形などは，領域を適切な変換で同一視して得る．例えば，位置だけを忘れる
なら平行移動同値類，位置と姿勢を忘れるなら合同変換同値類，尺度も忘れるなら相似同値類を
用いる．ただし，$M/G/E$のような反復商と形状商の重ね書きは，二つの作用の関係を指定しない
限り用いない．必要な変換で生成される同値関係を直接定義する．また，商の元は実現の型であり，
複数の実現そのものではない．

複合長パタンは，位置を失った同値類の単なる部分集合ではなく，順序や反復の担体を持つ
添字付き族として扱う．例えば「50分$\times$3コマ＋10分休憩」は，長さの集合ではなく
配列または周期付き列である．長さ値の空間自体には時空内の局所有限性を定義できないため，
必要な局所有限性はその実現またはアンカーの側に課す．

有界性・コンパクト性は$\mathcal{H}(M)$の定義には含めず，
$\operatorname{Bounded}(H)$，$\operatorname{Compact}(H)$のような独立の性質とする．
したがって，非有界な反復軌道も排除しない．

\subsection{時空間関連}

$A$側の分類は，時間と空間を分けず，次の二軸で行う．第一軸は点か領域か，第二軸は
単一か複合かであり，複合の下で繰り返しと恣意的な集まりを区別する．
\begin{figure}[H]
  \centering
  \begin{forest}
    dir tree
    [\tn{losktoli}{時空間関連}
      [\tn{toljoamili}{時空点パタン}
        [\tn{kon-toljoamili}{単一時空点}]
        [\tn{calf-toljoamili}{複合時空点}
          [\tn{flan-toljoamili}{繰り返し時空点}]
          [\tn{fazlo-toljoamili}{恣意的時空点}]
        ]
      ]
      [\tn{lokac-toljoafon}{時空領域パタン}
        [\tn{kon-toljoafon}{単一時空領域}]
        [\tn{calf-toljoafon}{複合時空領域}
          [\tn{flan-toljoafon}{繰り返し時空領域}]
          [\tn{fazlo-toljoafon}{恣意的時空領域}]
        ]
      ]
    ]
  \end{forest}
  \caption{参照構造から独立した時空間関連の分類}
\end{figure}

ここでflan-toljoamiliとflan-toljoafonは，時空間自身の対称性や，$M$上に直接定義された
生成規則による反復である．「毎日正午」のように暦・時計・タイムゾーンを必要とする反復は，
純粋な$A$側のパタンではなく，後述する$v$側の反復値である．

\subsection{時間値と空間値}

時間値・空間値は，参照構造$R$によって得られる時間・空間への分解を前提とする$v$側の
分類である．したがって，ここに時点・時区間・時間長を置くことは，時空間$M$そのものを
時間と空間へ絶対的に分割することを意味しない．

\begin{figure}[H]
  \centering
  \begin{forest}
    dir tree
    [\tn{tselselov}{時間値}
      [\tn{lokac-selovmili}{時点値パタン}
        [\tn{kon-selovmili}{単一時点値}]
        [\tn{calf-selovmili}{複合時点値}
          [\tn{flan-selovmili}{繰り返し時点値}]
          [\tn{fazlo-selovmili}{恣意的時点値}]
        ]
      ]
      [\tn{lokac-selovfon}{時区間値パタン}
        [\tn{kon-selovfon}{単一時区間値}]
        [\tn{calf-selovfon}{複合時区間値}
          [\tn{flan-selovfon}{繰り返し時区間値}]
          [\tn{fazlo-selovfon}{恣意的時区間値}]
        ]
      ]
      [\tn{lokac-selovtent}{時間長値パタン}
        [\tn{kon-selovtent}{単一時間長値}]
        [\tn{calf-selovtent}{複合時間長値}
          [\tn{flan-selovtent}{繰り返し時間長値}]
          [\tn{fazlo-selovtent}{恣意的時間長値}]
        ]
      ]
    ]
  \end{forest}
  \caption{時間値の分類}
\end{figure}

例えば，値の型と参照構造を明示すると，次のように分析できる．
\begin{itemize}
    \item 「今」は，発話文脈が供給する観測者と時点値を結ぶ指標的な値である．
    \item 「16:43」は，日付が別に与えられれば単一時点値になり，日内時刻の型として
    解釈すれば毎日反復する時点値パタンになる．
    \item 「2024年1月5日」は，暦とタイムゾーンのもとで一つの時空領域を指示する
    単一時区間値である．
    \item 「3時間」は，開始位置から独立した単一時間長値である．
    \item 「1か月」は，暦の規則によって実現長が変わりうる時間長記述である．
    \item 「毎日正午」は，参照構造のもとで複数の時空領域を生成する繰り返し時点値である．
\end{itemize}
以前の「時間的対象・位置記述・時間長・反復パタン」という四分は採用しない．この四者は
同一の分類軸上にないからである．本体系では，$A,R,v,h$の役割を先に分け，その後に$v$を
点・領域・長さで分け，さらに単一・複合・反復という生成様式を直交させる．

時点値・時区間値・時間長値の間の既存の関係は，$R$を共有する値どうしの関係として残せる．
例えば，単一時区間値の開始・終了・長さを次のように置く．
値水準の開始・終了関係は，生起物トークン間の境界関係\ko{moila}／\ko{fiila}（生起物の章，\ko{milistav}$\times$\ko{saistav}）と型が異なる．一語彙には一つの$\dom$しか与えられないので，共通の上位ドメインと制限定義を与えるまでは，値水準の関係を暫定記号$\mathrm{moila}_v$／$\mathrm{fiila}_v$として区別する（2026-09-03）．
\rele{moila}{kon-selovmili}{kon-selovfon}{（暫定記号$\mathrm{moila}_v$）$x_1$は$x_2$の開始時点値である}
\rele{fiila}{kon-selovmili}{kon-selovfon}{（暫定記号$\mathrm{fiila}_v$）$x_1$は$x_2$の終了時点値である}
\rele{cistoa}{kon-selovfon}{kon-selovtent}{$x_2$は$x_1$の時間長値である}
\rele{taks}{laccen}{kon-selovtent}{数値を所定の単位による時間長値へ対応させる}
これらは参照構造を省略した値水準の関係である．特定の生起物がいつ生じるかを述べる場合は，
生起物を時点値へ直接結ぶのではなく，その生起物の時空領域$A$について
$h=\langle A,R,v\rangle$を構成する．従来のakiは，この構成を文脈上省略した関係として扱う．

空間値も同じ構造を持つ．
\begin{figure}[H]
  \centering
  \begin{forest}
    dir tree
    [\tn{tselfenov}{空間値}
      [\tn{lokac-fenovmili}{場所値パタン}
        [\tn{kon-fenovmili}{単一場所値}]
        [\tn{calf-fenovmili}{複合場所値}
          [\tn{flan-fenovmili}{繰り返し場所値}]
          [\tn{fazlo-fenovmili}{恣意的場所値}]
        ]
      ]
      [\tn{lokac-fenant}{空間領域値パタン}
        [\tn{kon-fenant}{単一空間領域値}]
        [\tn{calf-fenant}{複合空間領域値}
          [\tn{flan-fenant}{繰り返し空間領域値}]
          [\tn{fazlo-fenant}{恣意的空間領域値}]
        ]
      ]
      [\tn{lokac-fenovtent}{広がり値パタン}
        [\tn{kon-fenovtent}{単一広がり値}]
        [\tn{calf-fenovtent}{複合広がり値}]
      ]
      [\tn{loskkelte}{方向値関連}
        [\tn{lokac-keltemili}{方向値パタン}
          [\tn{kon-keltemili}{単一方向値}]
          [\tn{calf-keltemili}{複合方向値}
            [\tn{flan-keltemili}{繰り返し方向値}]
            [\tn{fazlo-keltemili}{恣意的方向値}]
          ]
        ]
        [\tn{lokac-keltefon}{方向領域値}]
        [\tn{lokac-keltent}{方向広がり値パタン}]
      ]
    ]
  \end{forest}
  \caption{空間値の分類}
\end{figure}

方向は一般に$S^2$の元とは限らない．計量を備えた多様体では，点$p$における方向は
単位接球面$S(T_pM)$の元であり，場所をまたぐ方向場は単位接束の切断として表される．
「北」「右」「上」などの値は，参照構造が与える姿勢・重力・地磁気・対象固有軸などに
よって，この方向へ対応付けられる．なお，「未来向き」は空間方向ではなく，時空間の
因果構造が定める時間的向きであるため，空間方向値とは分ける．

\subsection{参照構造}
\todo{これをどこに置くのか悩ましい．暫定的にmetloskの下に置くのはあり．}
参照構造 notlileka は，単一の座標系に限らず，時空領域を値へ対応付けるために
必要な構成物を表す．一般形を
\begin{align*}
    R=(d,F,C,U,e,\gamma)
\end{align*}
と書く．ここで，$d$はアンカーまたはdatum，$F$は運動学的な基準構造，$C$は表示規則，
$U$は単位系，$e$はepochまたは零点，$\gamma$は粒度である．すべての位置付与が全成分を
必要とするわけではなく，不要な成分は省略できる．

「基準系」「観測者」「座標系」をis-aの木として並べることはしない．これらは指定の細かさが
異なる場合もあるが，存在論的には別の役割を持つからである．
\begin{itemize}
    \item 観測者は，典型的には一本の時間的世界線である．
    \item 基準系は，静止している観測者の世界線族を与える．必要なら同時性の切り方も伴う．
    \item 姿勢枠は，観測者に沿って空間方向の基底を与える．
    \item 数値座標系は，点に数値の組を割り当てる規則である．
    \item 暦・時計，住所・地名，順序尺度，intrinsic / relative / absoluteな質的参照枠は，
    数値座標系より広い表示規則$C$に属する．
\end{itemize}
例えば，グレゴリオ暦そのものは歴史的に制定された準抽象的な規約としてkika側に置けるが，
それをdatum・タイムゾーン・時法などと組み合わせて位置付与に用いる構造はnotlilekaである．
このように，参照構造を構成する各要素の存在論的分類と，複合体$R$が位置付与の基盤として
果たす役割は区別する．

葉層は基準系一般の同義語ではない．葉層とは，多様体を同じ次元の部分多様体へ局所的に
平行な形で分割したものであり，時空間では
\begin{align*}
    M=\bigcup_{t}\Sigma_t
\end{align*}
のような同時刻面の族が典型である．一方，静止観測者の世界線族はcongruenceまたはthreading
に相当する．例えばガリレイ時空$M=\R\times\R^3$で
\begin{align*}
    (t,\bm{x})\longmapsto
    (t+s,\,Q\bm{x}+\bm{v}t+\bm{a})
\end{align*}
と変換しても，$t=\text{const.}$という同時性葉層はすべての慣性系で共通である．速度
$\bm{v}$が変えるのは，むしろ
\begin{align*}
    u_{\bm{v}}
    =\partial_t+v^i\partial_i
\end{align*}
の積分曲線として与えられる世界線族である．これに対して特殊相対論では，観測者の
四元速度$u$に直交する超平面が同時性面となるため，速度の変更は同時性葉層も変更する．
回転系のように，大域的な同時性葉層を持たない基準系もある．

\subsection{時差と日内時刻}

時間表現にも参照構造への依存がある．「2024年1月5日」は暦とタイムゾーンに依存し，
「午前2時」は時法に依存し，「今」は発話文脈に依存する．一方，空間表現にも緯度経度の
ような絶対的参照枠と，「家の前」のような対象固有参照枠がある．したがって，時間を絶対的，
空間を相対的と存在論的に二分することはしない．ただし，自然言語では空間表現が物体基準の
参照枠を明示しやすく，暦日・時計時刻では規約的参照枠が省略されやすい，という使用上の
傾向は認められる．

時差の例では，普遍時間と日内時刻の位相を分ける必要がある．普遍時間を$t\in\R$，経度方向の
角度を$\theta\in S^1$とする．角速度$\omega$で回る座標への変換は
\begin{align*}
    \theta'&=\theta-\omega t,\\
    t'&=t
\end{align*}
であり，これはガリレイ的な時空座標変換である．一方，角速度$\omega$で経度方向へ移動する
観測者の世界線を$\theta(t)=\omega t+\theta_0$とすれば，その観測者に沿った日内時刻は
\begin{align*}
    \tau(t)
    =
    \left[(\omega_0-\omega)t+\tau_0\right]_{2\pi}
    \in S^1
\end{align*}
という派生ラベルで表される．ここで$\omega_0$は地球の自転角速度である．したがって，
\begin{align*}
    t_{\mathrm{local}}
    =
    \left(1-\frac{\omega}{\omega_0}\right)t
\end{align*}
という式は，位相を$\R$へ持ち上げて時間単位に換算した表示としては用いられるが，時空座標の
変換そのものではない．また，任意の実数倍$[t]\mapsto[ct]$は一般には$S^1\to S^1$の写像として
well-definedではないため，$t\in\R$から$\tau\in S^1$への写像として定義する．

\subsection{値の解釈と反復}

参照構造$R$のもとで値$v$が指示する時空領域の族を
\begin{align*}
    \llbracket v\rrbracket_R
    \subseteq
    \mathfrak{L}(M)
\end{align*}
と書く．単一値では$\llbracket v\rrbracket_R$が一つの点または領域を選び，複合値では複数の
点・領域を選ぶ．canontoli $h=\langle A,R,v\rangle$が成立するための基本条件は
\begin{align*}
    A\in\llbracket v\rrbracket_R
\end{align*}
である．値が曖昧な場合には，右辺は複数の候補を含む．精度や曖昧性を表す必要がある場合は，
$v$の型または別の記述対象にその情報を持たせる．

反復には，次の二種類がある．
\begin{enumerate}
    \item $A$側の内在的反復：$M$上の群作用や生成規則から直接得られる$G\cdot A_0$，
    \item $v$側の参照相対的反復：暦・時計・格子座標などを用いて得られる
    $\llbracket v\rrbracket_R=\{A_i\mid i\in I\}$．
\end{enumerate}
例えば，結晶格子の対称性を$M$上に直接与えれば前者として扱え，「毎日正午」は後者として
扱う．同じ外延を持つ場合でも，生成規則と同一性基準が異なるなら別のパタンとして区別する．

以上により，本節の基本構造は
\begin{align*}
    \boxed{
      \text{時空領域 }A
      \;--\;
      \text{参照構造 }R
      \;--\;
      \text{未束縛の値 }v
      \;--\;
      \text{参照付き位置 }h
    }
\end{align*}
とまとめられる．時間と空間の分離，点・領域・長さの区別，単一・複合・反復の区別は，
この基本構造より後の，互いに異なる分類軸として導入される．

\subsection{kika}
\ko{kika}は，従来は主として表現を扱う準抽象物の概念クラスとされてきた．しかし，表現，表現形態，内容，物理的な表現実現物，虚構世界内の対象が十分に区別されておらず，抽象物\ko{loka}との境界も未確定である．この問題群は独立した第\ref{chap:kika}章で扱う．現段階では旧分類木を採択済みの木として掲げない．
\section{形式的対象の分類}\label{sec:stcilkant-classification}

\ko{stcilkant}は，同一性基準が物理世界への参照なしに形式体系の内部で
完結する対象である．この特徴づけは形式的対象の実在を主張するものではなく，
何を同じ対象と数えるかを与える作業的な分類基準である．数を特定の集合と同一視せず，
構造の位置として扱う方針は，同型な集合論的実装が複数あることを指摘した
\citet{benacerraf1965}以降の数学的構造主義と両立する\citep{shapiro1997}．

\ko{stcilkant}の直下は，同一性の基準が違う形式的対象の種類を外延的に並べたコンストラクタ分割（列挙）とし，兄弟は排他とする（2026-09-30）．数の成員は個々の数であり，構造としての$\mathbb Z$（環）とは別の対象である．平面の中の円（図形）と微分同相で同一視する$S^1$（多様体）も同一性の基準が違う．数・写像・図形を集合の下に置かないのは，数を特定の集合と同一視しない方針による．重なるのは構造どうしであり（束は順序構造でも代数構造でもあり，位相群は空間でも群でもある），〈構造〉の下の非排他の素性分割とする．論理式・命題は形態の記号列と重なるので上位未決に置く．網羅的でなく，後述の構造スキーマや関係グラフの代わりでもない．〈〉は未造語．
\begin{figure}[H]
  \centering
  \begin{forest}
    dir tree
    [stcilkant\\形式的対象
      [laccen\\数
        [mistakk\\自然数]
        [stakk\\整数]
        [flakk\\有理数]
        [tolas\\実数]
        [klanlas\\複素数]
      ]
      [fonne\\集合
        [〈形式言語〉]
      ]
      [kilemhaisk\\形態
        [〈記号〉]
        [suyflos\\記号列]
        [〈構文木〉]
      ]
      [mintos\\論理式・命題]
      [haisk\\図形]
      [koln\\写像・射]
      [〈構造〉
        [〈代数構造〉
          [〈群〉]
          [〈環〉]
          [kolnfin\\圏]
        ]
        [〈空間〉
          [〈位相空間〉
            [saljaaflemae\\多様体]
          ]
          [〈可測空間〉]
        ]
        [〈順序構造〉]
      ]
      [〈理論〉]
      [〈論理体系〉]
    ]
  \end{forest}
  \caption{\ko{stcilkant}の分類木（2026-09-30，暫定）}
\end{figure}

\paragraph{形態と記号的対象} 記号・記号列・構文木は，表現章が\ko{stcilkant}に置く形態\ko{kilemhaisk}の下に置き，コンストラクタ分割を立てる：記号（原子），記号列（連接），構文木（構成規則による組み立て）．ここでの記号は形式体系のアルファベットの元であり，文字や矢印のような記号の型（\ko{suylis}）は制作と共同体に依存する\ko{kika}，その書かれた実現は\ko{haikant}である．三者は同一性の基準が違う．記号はアルファベットの中で，記号列は列として，構文木は木として同じなら同じであり，同じ記号列が別の構文木を持ちうる．旋律や線画のように記号的でない形態はこの分割の外にある（2026-09-29）．次のものは記号的対象でなく，別の場所に置く．
\begin{itemize}
  \item \emph{形式言語}：記号列の集合．\ko{fonne}の一種である．
  \item \emph{論理体系}：言語と帰結関係の組．次節の仕様層の側にあり，個々の理論（シグネチャと公理系の組）とも区別する．
  \item \emph{公理・定理}：理論を文脈とするロール（\ko{liflom}）．担い手は論理式だが，論理式そのものではない．「冪集合の公理」はZFCを文脈として公理ロールを担う論理式の名である．
  \item \emph{コノメノの表現}：形式文法章の表現の埋め込み$\ulcorner\cdot\urcorner$の像としての構文は記号列・構文木だが，規約の束としての言語や個々の発話は\ko{kika}・\ko{astav}である（表現章「言語の立ち位置」）．語義を分けて書く．
\end{itemize}

\subsection{理論，表示，抽象構造}

代数系や位相空間を単に「順序組である」と定義すると，順序組そのものの集合論的実装に
同一性が依存してしまう．そこで，次の三層を分ける．
\begin{enumerate}
  \item \textbf{仕様層}：ソート，演算記号，関係記号からなるシグネチャ$\Sigma$と，
  公理系$T$．$T$は構造そのものではなく，許される構造を特徴づける形式的記述である．
  \item \textbf{表示層}：集合論などの基礎体系内で構成した$\Sigma$構造$\mathcal A$．
  これは$\mathcal A\models T$を満たす一つのモデルである．
  \item \textbf{抽象構造層}：表示が指示する数学的対象．語義ごとに指定した同型または
  それより弱い同値関係$\simeq_E$に対して不変なものとして扱う．
\end{enumerate}
この三層は，数学的概念を言語非依存のドメイン層と言語層に分ける
OntoMath$^{\mathrm{PRO}}$ 2.0の方針\citep{kirillovich-etal2022}や，数学的対象の抽象表現と
アプリケーション内部の表現，通信用符号化を分けるOpenMathの方針\citep{openmath2017}に対応する．
ただし，これらの実装上の層をそのまま形而上学的な答えとはしない．

多ソートのシグネチャ$\Sigma=(S,F,R)$に対する表示は，例えば
\begin{align*}
  \mathcal A
  =
  \left\langle
    (A_s)_{s\in S},
    (f^{\mathcal A})_{f\in F},
    (r^{\mathcal A})_{r\in R}
  \right\rangle
\end{align*}
と書ける．群の表示$\langle G,\cdot,e,{}^{-1}\rangle$は群の公理を満たす代数構造，
位相空間の表示$\langle X,\tau\rangle$は$\tau\subseteq\mathcal P(X)$が位相の公理を満たす構造である．
多様体，順序構造，グラフ，圏も，必要なシグネチャと公理を変えて同じ枠組みに載せる．
ただし，一般の数学実践では同一視の基準が常に同型とは限らず，圏の同値，ホモトピー同値などを
用いる領域もある．そのため，\ko{stcilkant}全体に一つの同型関係を固定せず，
概念ごとに$E$を記録する．

\subsection{下位分類，埋め込み，一般化}

数学的対象の間には，少なくとも次の異なる関係がある．
\begin{center}
\begin{tabular}{p{0.22\linewidth}p{0.68\linewidth}}
\toprule
関係 & 読み方 \\
\midrule
$A\sqsubseteq B$ & 概念クラスの下位分類．$A$のすべての個体が$B$の個体である． \\
$a\in A$ & 対象と集合または数学的コレクションの成員関係． \\
$A\hookrightarrow B$ & 構造を保つ埋め込み．端点を同一視するかどうかは別に指定する． \\
$A\simeq_E B$ & 指定した構造に関する同型または同値． \\
$U:\mathcal A\mapsto\mathcal A|_{\Sigma'}$ & 演算や関係の一部を忘れる忘却・縮約．構造の一般化に用いる． \\
\bottomrule
\end{tabular}
\end{center}

この区別により，自然数，整数，有理数，実数，複素数は\ko{laccen}の下の兄弟分類とする．is-aで上へ一般化しないのは，拡張が枝分かれするからである：$\mathbb Q$は$\mathbb R$にも$p$進数$\mathbb Q_p$にも拡がり，$\mathbb N$は$\mathbb Z$にも順序数にも拡がる．包含の鎖として積むと，数は両立しない拡張の和になって木にならず，頂点も決まらない（四元数・八元数の先も続く）．一般化は次の埋め込みのグラフで持つ（2026-09-30）．
標準的な関係は
\begin{align*}
  \mathbb N\hookrightarrow\mathbb Z\hookrightarrow\mathbb Q
  \hookrightarrow\mathbb R\hookrightarrow\mathbb C
\end{align*}
と書く．$\mathbb N\subseteq\mathbb C$は，これらの埋め込みによって元を同一視する特定の実装規約を
採った後なら使ってよいが，\ko{skon}関係そのものとはしない．

普通の分類木は，上位概念の外延を下位概念が分割するため下側に分岐する．これに対し，
一般化は「どの構造を忘れるか」による関係なので，一つの具体的な構造から複数の一般化先が上側へ分岐しうる．
例えば環は，加法構造を取れば可換群，乗法構造を取ればモノイドへ映る．これは分類木の向きが
反転したのではなく，異なる関係を見ているためである．理論の再利用には，単純な多重継承ではなく，
記号の対応を指定する理論射のグラフを用いる方がよい\citep{omdoc12}．

したがって，\ko{stcilkant}には一本の排他的なskon木だけを持たせない．
\begin{enumerate}
  \item コンストラクタ分割は，同一性の基準が違う種類（数，集合，図形，形態，写像，構造など）の外延的な列挙に限る．構造どうしの重なりは〈構造〉の下の素性分割で書く．
  \item 構造付き対象は$\langle\Sigma,T,\mathcal A,E\rangle$を参照する形式で記述する．
  \item 埋め込み，同型・同値，忘却，理論射を，skonとは別のラベル付き関係として保存する．
\end{enumerate}

\begin{remark}[公理系は対象の実装ではない]
自然数や図形を，それを表示する集合や特定の公理系から分ける方針は維持する．
ただし，公理系$T$はモデルのクラスを特徴づけるもので，集合論的構成が一つのモデルを与えることとは異なる．
特に，一階のPeano算術のように非標準モデルを持つ理論は，公理だけで意図した一つの構造をカテゴリカルに
決定しない．また，論理式の記号列，その構文木，命題内容は異なる同一性基準を持つため，
\ko{mintos}が「論理式」と「命題」を兼ねる現行訳は暫定とする．
\end{remark}

\section{現実対象側の抽象物}

\subsection{loka}

\ko{loka}は\ko{suiskant}の下で，時空間に局在しないが，現実世界の質，測定，
記述実践などに関与する抽象物を扱う．現行の主な下位分類は次の通りである．
\begin{figure}[H]
  \centering
  \begin{forest}
    dir tree
    [\tn{loka}{抽象物}
      [\tn{latent}{量}
        [\tn{lekostent}{定値量}]
        [\tn{nosnant}{形容量}]
      ]
      [\tn{liflom}{ロール}]
      [(評価記述)]
      [(傾向タイプ)]
      [(形式構築物)]
      [(概念個体)]
    ]
  \end{forest}
  \caption{\ko{loka}の現行の概略}
\end{figure}

\paragraph{概念個体（未造語）} 辞書項目そのもの——分類木の節点，量種（長さ・配位・所有者），行為ロール概念——を個体として\ko{loka}下に置く（2026-09-02）．旧版のhalko（品詞の品詞）が既に前提していた扱いの明文化である．必要性は二箇所から来る：
\begin{enumerate*}
    \item 単純生起物の射影$q$は量種であり（\ko{latent}の個体は18cmのような値なので，量種は値のクラスに当たる），これを一階の量化域に入れるため，
    \item 素性による定義クラス（「殺す」＝「死ぬ」$\land\exists$\ko{towa}）や$\mathrm{Tel}$（完了条件つき概念，生起物の章）が概念の上の述語であるため．
\end{enumerate*}
種を立てる基準では4（空外延・共外延）と6（同一性の様式の供給）に当たる．同一性の様式は二種を宣言する：原生概念は原始的（開ハブ），構成された概念（$\qquote{Rn}$等）は成分上の閉ハブ（形式的構築による閉ハブの判例）．表現（\ko{kika}）—指示→概念個体—$\rho_w$→外延という三層で，$\qquote{\cdot}$は第一層から第二層への写像の記法，フレームの語彙割り当て$\rho_w$は概念個体から外延への割り当てと読み替える（形式意味論の章の定義は変えずに済む）．概念個体のクラス自身も概念個体である（再帰と高階の許容）．
\ques{名称，および\ko{latent}・\ko{liflom}との位置関係（量種は概念個体の特殊例か）．}

\end{document}
